\subsection*{集合论的悖论与基础危机}

第一批“悖论性”集合出现在基数和序数的理论中。1897年，布拉利-福尔蒂注意到，不可能存在由所有序数构成的集合，因为这样的集合将是良序的，因此与自身的某个真段同构，这是荒谬的（*）。1899年，康托尔（在给戴德金的信中）观察到，我们既不能再说基数构成一个集合，也不能谈论“所有集合的集合”而不产生矛盾（后者“集合”Ω的所有子集所构成的集合将与Ω的某个子集等势，这与不等式m \textless{} 2 \(^{m}\) ) ({[}25{]}，第444-448页）。最后，在1905年，罗素通过分析这一不等式的证明表明，确立该不等式的推理同时也证明了（无需诉诸基数理论）“所有不包含自身作为元素的集合的集合”这一概念是自相矛盾的 {[}37{]} (†).

人们可能会认为，这样的“悖论”只会出现在数学的边缘领域，其特征是考虑那些“大”到直觉无法企及的集合。但其他悖论很快就威胁到了数学中最经典的部分。贝里和罗素 {[}37{]}简化了J. Richard的一个论证， {[}34{]}观察到可以用少于十九个音节来命名的整数集合是有限的，但尽管如此，将整数定义为“不能用少于十九个音节命名的最小整数”却是矛盾的，因为这个定义只包含十八个音节。

尽管这类论证远离数学家当前的用法，在许多人心目中必定显得像是一种文字游戏，但它们却表明有必要修订数学的基础，以消除这类“悖论”。然而，尽管对于这种修订的紧迫性意见一致，但在如何实现修订的方式上很快就出现了根本性的分歧。对于一组数学家，即“唯心主义者”和“形式主义者”（*），集合论“悖论”所造成的情况与非欧几何或“病态”曲线（如无切线的曲线）的发现所引起的情况非常相似；因此，它应当导致一个类似但更一般的结论，即试图将任何数学理论建立在诉诸（明示或其他方式的）“直觉”之上是徒劳的。这一立场可以用形式主义学派主要对手的话来概括：“\ldots{} 形式主义的观点”，布劳威尔说（{[}38 a{]}, 第83页），“认为人类理性既没有关于直线也没有关于大于十的数的精确图像，例如，\ldots{} 诚然，从我们假定为公理的数学实体之间的某些关系中，我们按照固定法则推导出其他关系，并确信通过这种方式我们借助逻辑推理从真理得出真理。\ldots{} 因此，对于形式主义者来说，数学的精确性仅仅在于发展关系序列的方法，而与人们可能想赋予这些关系或其关联实体的意义无关。”

因此，对于形式主义者来说，所需要的是为集合论提供一个与初等几何的公理化基础完全类似的基础，在其中既不必知道被称为“集合”的“事物”是什么，也不必知道关系 \(x \in y\) 意味着，但其中列举了对这一关系所施加的条件；自然，这样做的方式应当尽可能涵盖康托尔理论的所有成果，同时使得“悖论性”集合的存在成为不可能。这种公理化的第一个例子由策梅洛于1908年给出。 {[}36 c{]}他避免了“过大”的集合，通过引入一个“分离公理”（Aussonderungsaxiom），根据该公理，一个性质 \(P\{x\}\) 确定一个由具有此性质的元素构成的集合，仅当 \(P\{x\}\) 已经蕴含了如下形式的关系 \(x \in A (\dagger)\) 。但只有通过限制“性质”这一概念所附带的含义，才能消除类似于“理查德悖论”的悖论。

在这方面，策梅洛满足于用极其模糊的措辞来描述一类他称之为“确定的”性质，并指出选择公理的应用必须限于这类性质。这一点由斯科伦 {[}39{]} 和弗伦克尔 {[}41{]} (*) 加以精确化；正如他们所观察到的，对其阐明需要一个完全形式化的系统（如本书第一、二章所描述的那样），在该系统中，“性质”和“关系”的概念已失去一切“意义”，而仅仅成为按照明确法则形成的组合体的名称。这当然要求所使用的逻辑规则被纳入该系统，而这在策梅洛-弗伦克尔系统中尚非如此；除此之外，我们在第一、二章中所描述的本质上正是他们的系统。

随后，人们提出了集合论的其他公理化方案。我们将主要考虑冯·诺依曼的体系 {[}43 a 和 b{]}，它比策梅洛-弗伦克尔体系更接近康托尔的原始构想；后者为了避免“悖论性”集合，已经（在与戴德金的通信中， {[}25{]}第445-448页）提议区分两类集合，即“多重性”（Vielheiten）和严格意义上的“集合”（Mengen），后者以能够被思考为单一对象为特征。正是这一思想由冯·诺伊曼通过区分两类对象——“集合”和“类”——而加以精确化。在他的体系（几乎完全形式化）中，类与集合的区别在于它们不能置于符号 \(\in\) 的左侧。这类体系的一个优点是它恢复了十九世纪逻辑学家所使用的“全类”概念（当然，它不是集合）的地位。我们还应当补充说明，冯·诺伊曼的体系避免了（对于集合论而言）引入公理模式，这些公理模式被适当的公理所取代，从而使该体系的逻辑研究更加容易。冯·诺伊曼体系的变体由贝尔奈斯和哥德尔给出 {[}44 b{]}.

这些体系似乎已经成功地消除了悖论，但代价是那些不能不显得极其任意的限制。为策梅洛-弗伦克尔体系辩护可以说，它仅限于制定一些禁令，而这些禁令不过是对集合概念在各种数学理论应用中现行做法的认可。冯·诺伊曼和哥德尔的体系则与通常的观念相距较远。另一方面，我们不能排除这样一种可能性，即把某些数学理论的基础纳入这类体系的框架，可能比纳入策梅洛-弗伦克尔体系更为僵硬的框架更容易。

无论如何，不能说这些解决方案中的任何一个给人以最终定论的感觉。它们满足了形式主义者的要求，因为形式主义者拒绝考虑每位数学家的个体心理反应；他们认为，一种形式化语言如果能够以不含歧义的形式转写数学论证，从而充当数学思想的载体，那么它就尽到了自己的职责。他们会说，每个人都有权对他们所使用数学实体的“本性”或定理的“真理性”持有自己的看法，只要他们的推理能够被转写为共同语言（*）。

换言之，从哲学观点来看，形式主义者的态度特征在于对“悖论”所提出的问题漠不关心，这种态度源于放弃了柏拉图式的立场，即主张赋予数学概念以所有数学家共有的理智“内容”。许多数学家反抗这种与传统的决裂。例如，罗素试图通过更深入地分析悖论的结构来避免它们。罗素和怀特海采纳了J.理查德（在文章 {[}34{]} 中，他提出了他的“悖论”）首次提出、后来由H.庞加莱 {[}33 c{]}发展的一个思想，他们观察到，悖论性集合的定义都违反了以下被称为“恶性循环原则”的原理：“凡涉及一个集合的全部者，必不属于该集合”（{[}37{]}，第I卷，第40页）。这一陈述成为《数学原理》的基础，并且为了容纳它而发展了“类型论”。与启发它的弗雷格逻辑一样，罗素和怀特海的逻辑（远比本书第一章中使用的数理逻辑更为宏大）包含“命题变项”。类型论对这些变项进行分类，其大致轮廓如下。

从未定义的“个体域”出发，该域可被称为“0阶对象”，其中变量（自由或约束）为个体的关系被称为“一阶对象”；一般而言，变量为阶 \(\leqslant n\) 的对象（且其中至少有一个是 n 阶的）的那些关系，被称为“ \(n + 1\) 阶对象”（*）。因此，一个 n 阶对象的集合只能通过一个 \(n + 1\) 阶关系来定义，且此条件使得悖论集合的消除毫无困难（ \(\dagger\) ).但“类型层级”原则的限制性如此之强，若严格遵循，将导致数学陷入无法解脱的复杂性（**）。为摆脱这一后果，罗素和怀特海不得不引入一条“可化归性公理”，该公理断言，对于“个体”之间的每一种关系，都存在一个等价的一阶关系。这一条件与形式主义者的公理一样武断，并大大降低了《数学原理》建构的兴趣。罗素和怀特海的体系在逻辑学家中比在数学家中更受欢迎。此外，它并未完全形式化（ \(\dagger\dagger\) )，因此在细节上存在许多模糊之处。人们曾多次尝试简化和澄清它（拉姆齐、赫维斯特克、奎因、罗瑟）；这些作者倾向于使用越来越完全形式化的语言，并将《数学原理》中的规则（这些规则仍含有某种直观基础）替换为仅考虑所讨论的组装书写方式的限制；这些规则不仅因此显得像策梅洛-弗伦克尔或冯·诺依曼体系中规定的禁令一样武断，而且由于更远离数学实践，在许多情况下导致了作者未曾预见的不可接受后果（如布拉利-福尔蒂悖论，或选择公理的否定）。

对于这些学派的数学家而言，根本目标在于不放弃过去遗产的任何部分。希尔伯特曾说：“没有人能把我们赶出去。”{[}30{]}，第274页），“从康托尔为我们创造的天堂”。为了实现这一目标，他们准备接受数学推理上的限制，这些限制虽然因符合数学惯例而不具有至关重要性，但似乎并非由我们的思维习惯和集合的直觉概念所强加。对他们而言，任何事物都比心理学侵入数学有效性标准更可取；与其像阿达马所说的那样“考虑我们大脑的性质”（{[}35{]}，第270页），他们宁愿接受大体上任意划定的边界，只要所有经典数学都包含其中，并且这些边界不会威胁到阻碍未来的进展。

截然不同的是坚持下述学说的数学家们的态度。形式主义者同意在数学推理方面放弃“心灵之眼”的控制，而被贴上“经验主义者”、“实在论者”或“直觉主义者”标签的数学家则拒绝同意这种放弃；他们坚持某种内在的确定性，以保证他们所研究的数学对象的“存在”。他们对放弃空间直觉没有严重异议，因为算术“模型”使他们能够躲藏在整数的直觉概念之后。但在将整数概念还原为集合概念（这在直觉上远不清晰），然后为集合的操纵设置缺乏直觉基础的障碍这一问题上，产生了不可调和的反对。第一个表达这种反对的人（也是由于天才的权威而最具影响力的人）是H.庞加莱。他不仅接受了关于几何的公理观点和分析的算术化，还接受了康托尔理论的很大一部分，但他拒绝认为算术也可能适合公理化处理。特别是数学归纳原理，在他眼中是我们智力的基本直觉，不能纯粹被视为一种约定（*） {[}33 c{]}。他原则上反对形式化语言，质疑其有用性，并且倾向于将形式化数学中的整数概念与证明论（当时正在兴起，我们稍后将讨论）中对整数的使用相混淆。毫无疑问，在当时要像我们今天这样精确地做出这一区分并不容易——经过五十多年的研究和讨论之后——尽管希尔伯特和罗素等人已经清楚地把握了这一区分。

在1904年策梅洛引入选择公理之后，这类批评成倍增加 {[}36 a{]}。此前在分析和集合论的许多证明中，它的使用几乎未被注意到（*），正是沿着埃哈德·施密特提出的一个想法，策梅洛明确陈述了这一公理（并以《结果概要》§4第10号中给出的两种形式），通过一种巧妙的方法（我们在第三章§2中实质上复现了该方法）推导出了任何集合上良序存在的令人满意的证明。看来，由于它与“悖论”同时出现，这种新颖的推理方法以其陌生的面貌使许多数学家感到困惑；只需阅读《数学年鉴》下一卷中那些对康托尔方法十分熟悉的作者（如舍恩弗利斯和F.伯恩斯坦）由此产生的奇怪误解即可。E.博雷尔在同一卷中发表的批评更有实质内容，清晰地反映了庞加莱关于整数的观点；这些批评在E.博雷尔、贝尔、阿达马和勒贝格之间的书信往来中得到了发展和讨论，这已成为法国数学传统中的经典文献 {[}35{]}。博雷尔首先否认选择公理的有效性，因为一般而言它涉及不可数无穷多次选择，这在直觉上是无法想象的。但阿达马和勒贝格指出，可数无穷多次连续任意选择同样不具有直觉性，因为它涉及无穷多次操作，而这些操作不可能被认为实际完成。在扩大了讨论范围的勒贝格看来，一切都取决于“数学对象‘存在’”这一断言的含义；对他而言，必须明确“命名”一个性质（我们应称之为“泛函”性质），该性质唯一地定义该对象。像策梅洛在其证明中所用的那种函数，正是勒贝格所称的“选择法则”。他接着说，如果这一要求得不到满足，而我们只是“考虑”这个函数而非“命名”它，那么我们能否确信在整个推理过程中我们始终考虑的是同一个函数？{[}35{]}（第267页）这一考虑使勒贝格产生了新的疑虑：即使是在一个集合中选择单个元素似乎也会带来困难；必须确信这样的元素“存在”，即该集合中至少有一个元素可以被“命名”（†）。因此，如果我们不能“命名”一个集合的每一个元素，我们还能谈论该集合的“存在”吗？贝尔毫不犹豫地否认了给定无限集的子集之集的“存在”（同上，第263-264页）；阿达马徒劳地指出，这些限制将导致甚至放弃谈论实数集的权利，而E.博雷尔最终也得出了这一结论。除了可数集似乎被赋予了选举权这一事实之外，我们几乎又回到了“实无限”反对者的经典立场。

这些反对意见中没有哪一条是非常系统的；将数学按照类似但更为激进的路线进行全面重构的任务留给了布劳威尔及其学派。我们在此不打算贸然概括像直觉主义这样复杂的学说——它既是数学也是哲学——我们只打算指出它的一些最显著的特征，更详细的内容请参阅布劳威尔本人的著作 {[}38{]} 以及海廷的阐述 {[}46{]} 。在布劳威尔看来，数学与我们思想中“精确”的部分是同一的，这一部分奠基于自然数序列的原初直觉，并且若不加以割裂，就无法被转译成形式系统。此外，它只在数学家的心智中才是“精确”的，而期望构建一种完全摆脱语言一切缺陷与歧义的数学家之间的交流工具，乃是一种幻觉；最多只能期望在对话者心中唤起一种有利于或多或少模糊描述的心态（{[}46{]}，第11-13页）。直觉主义数学赋予逻辑的重要性几乎不高于语言；一个证明之所以具有决定性，并非凭借一经确定便永不变更的逻辑规则，而是因为其每一步都具有“直接的自明性”。此外，这种“自明性”还须以比E.博雷尔及其支持者所理解的意义更为严格的方式来解释；因此，在直觉主义数学中，除非对于出现在R中的变量所取的每一个值系统，我们都能证明命题R与“非R”二者之一为真，否则我们不得断言形如“R或（非R）”的关系为真（排中律）。例如，由方程 \(ab = 0\) 在两个实数之间，我们不能推断“ \(a = 0\) 要么 \(b = 0\) ``，因为很容易构造实数的显式例子 \(a\) , \(b\) 使得我们有 \(ab = 0\) 目前尚无法证明这两个命题中的任何一个 \(a = 0\) , \(b = 0\) ({[}36{]}，第21页）。

直觉主义数学家从这样的原则出发，得出与经典定理截然不同的结果，这并不令人意外。后者中有大量定理消失了，例如分析学中大多数“存在性”定理（如关于实函数的博尔扎诺和魏尔斯特拉斯定理）；如果一个实变量函数在直觉主义意义上“存在”，它便理所当然地是连续的，而一个单调有界的实数序列不一定有极限。此外，许多经典概念在直觉主义数学中分化为几个根本不同的概念。因此，收敛（对于实数序列）有两种概念，可数性有八种概念。不言而喻，超限归纳法及其在现代分析中的应用，如同康托尔理论的大部分内容一样，被毫无上诉余地地否定了。

根据布劳威尔的观点，只有通过这种方式，数学命题才能获得“内容”；超出直觉主义所承认范围的形式主义论证被认为毫无价值，因为无法赋予它们一种“意义”，使得直觉上的“真”概念能够适用。显然，这类判断只能基于先前对“真”的概念，而这一概念具有心理或形而上学的性质；也就是说，在实践中，它们超出了讨论的范围。

当然，来自直觉主义阵营的猛烈攻击不时迫使不仅前卫数学学派，甚至传统数学的拥护者，都不得不采取守势。一位著名数学家曾承认，他对这些攻击印象深刻，以至于自愿将自己的工作限制在那些被认为是“确定无疑”的数学分支中。但此类情况想必相当罕见。直觉主义学派，其记忆无疑将仅作为历史奇闻留存，至少提供了一项服务，即迫使它的反对者，也就是绝大多数数学家，澄清各自的立场，并更自觉地意识到他们对数学抱有信心的原因（无论是逻辑上的还是情感上的）。

